\section*{Ex.\ 1.0 --- Three Papers on OS History and Architecture}

The three papers below were located by searching the literature (Google Scholar and
general web search for work on operating system history and architecture), downloaded in
full and read; none of them is taken from the supplied paper list. They were chosen as a
set because they form a single argument across forty-four years. Each gives a different
answer to the same two questions: \emph{where should operating system function live}, and
\emph{what may the system assume about the machine underneath it?} Multics answers with one
comprehensive system on a symmetric shared-memory machine; Mach answers by decomposing that
system into a minimal kernel plus user-level servers; the Multikernel answers by discarding
shared memory as the organising assumption altogether. Read in order, they also show a
steady rise in evidentiary standards --- from a design prospectus with no measurements at
all, to a status report resting on a single microbenchmark, to a broad evaluation that
volunteers its own weaknesses --- without any matching fall in the confidence of the
architectural claim being made.

\subsection*{Paper 1 --- Corbat\'o and Vyssotsky, ``Introduction and Overview of the Multics System'' (1965)}
\label{p4:sec:multics}

\textbf{Summary.} The paper opens the six-paper Multics session at the 1965 Fall Joint
Computer Conference and states the objectives the other five then elaborate. Its problem is
access: batch processing had, in the authors' words, caused user access to large machines to
``retrogress,'' isolating the programmer from cause and effect. Their goal is a
\emph{computer utility} --- a system running continuously and reliably twenty-four hours a
day like a telephone or power network, serving everything from interactive terminals to
absentee (batch) jobs. The mechanisms proposed are largely hardware--software co-design on
the GE~645, a modification of the GE~635. Addressing is two-dimensional: programs are
written as \emph{segments}, each of which may grow or shrink during execution (up to roughly
a quarter of a million segments per user, each of up to a quarter of a million words), and
segments are independently \emph{paged} by the supervisor with a choice of 64- or 1024-word
pages. Cross-segment references are bound by \emph{dynamic linking} --- a segment names
another only symbolically, binding occurs automatically on first execution, and afterwards
runs at full speed. Pure, recursive, shareable procedures are the normal mode of
programming, so that one copy of a procedure serves many users and, as the authors note,
there is ``no clear-cut demarcation'' between user and system programs. Protection comes
from descriptor bits acting as hardware ``fire-walls,'' including an execute-only bit, and
most supervisor modules run with the same descriptors as user code and no access to
privileged instructions. Processors form an anonymous pool with no master --- ``the
supervisor does not have a special processor'' --- and reach I/O only through mailboxes in
the memory modules. A file system names everything symbolically rather than by address,
interlocks writing, and provides system-run backup against three named contingencies. I/O
is device-independent, batch is simply the zero-terminal case of an ordinary process,
resource budgets are hierarchical and delegable, PL/I is adopted for machine independence,
and the reference manual is kept on-line. There is no evaluation: the system did not yet
exist. The sole quantitative claim is that ``a few hundred'' simultaneous users are expected,
derived by ``simple scaling of processor and memory speed'' and immediately qualified as
``hazardous.'' The paper closes by calling the plans ``not unattainable'' while conceding it
would be ``presumptuous'' to think the initial system could meet all the stated requirements.

\textbf{Critical judgment.} This is a prospectus, and it cannot be judged as an empirical
paper: every verb is in the future tense and nothing in it is falsifiable by its own
evidence. The deeper omission is not measurement but a \emph{cost} estimate. The paper
proposes to meet ``almost all of the present and near-future requirements'' of a large
installation and never once asks what such generality would cost to build. That is precisely
where the project failed --- development ran years late, Bell Labs withdrew in 1969, and
UNIX was written shortly afterwards as a deliberate reaction against exactly this scale.
The most accurate sentences in the paper are its own hedges. Elsewhere it argues
rhetorically rather than evidentially: ``machine-aided cognition,'' the promised social
consequences, and above all the claim that PL/I would deliver machine independence are
asserted rather than demonstrated, and the last of these badly underestimated the cost of
building the compiler. On substance, history has split the paper cleanly in two. Vindicated:
paged virtual memory, device-independent I/O, a symbolically named hierarchical file system
with per-file access control and system-provided backup, symmetric multiprocessing with no
master processor, writing an operating system in a high-level language, and graded hardware
protection --- the descriptor bits became Multics' rings, which x86 inherited. Refuted:
user-visible \emph{segmentation} lost decisively to the flat, paged address space, and
x86-64 all but abolished segmentation; nobody today programs a two-dimensional address
space, and although dynamic linking survived it did so as lazy load-time binding over a flat
space, not as segment binding. The most interesting case is the computer utility itself.
The economic argument --- that bulk economies, floor space and management efficiency favour
one maximal centralised machine --- was falsified for roughly thirty years by the
minicomputer and the personal computer, then vindicated by cloud computing. The paper was
right about the destination and wrong about both the timing and the reason, since what
eventually made centralisation win was statistical multiplexing at datacentre scale, not the
considerations it actually advanced.

\subsection*{Paper 2 --- Accetta et al., ``Mach: A New Kernel Foundation for UNIX Development'' (1986)}
\label{p4:sec:mach}

\textbf{Summary.} Mach's target is the accretion of mechanism in UNIX. The authors observe
that a system which began with a small uniform set of operations on file descriptors had
grown ``a staggering number'' of overlapping facilities --- System~V streams, 4.2BSD
sockets, ptys, several kinds of semaphore, shared memory, and a ``mind-boggling array'' of
\texttt{ioctl} operations --- giving non-uniform access to resources both within a machine
and across a network, while the process abstraction remained too heavyweight to express
parallelism. Mach replaces the base with four abstractions. A \emph{task} is the unit of
resource allocation, holding a paged address space and a set of port rights; a \emph{thread}
is the unit of CPU utilisation, so that a UNIX process is simply a task with one thread; a
\emph{port} is a kernel-protected message queue acting as a capability, with any number of
senders but exactly one receiver; and a \emph{message} is a typed collection of data that may
carry port rights and may be as large as an entire address space. Every operation on any
other object is a message to the port representing it, so that the kernel is merely one
server among many --- the design thesis is that ``the actual system running on any particular
machine is a function of its servers rather than its kernel.'' The virtual memory system
supports per-page inheritance (shared, copy, or none) and a current-versus-maximum protection
pair in which the maximum may only ever be lowered; copy-on-write implements both
\texttt{fork} and bulk message transfer, illustrated with a 24\,MB message that passes
through a temporary kernel map. Most consequentially, page faults may be serviced
\emph{outside} the kernel by user-level pagers, so that a memory-mapped file is simply memory
whose pager is the file system, and the default pager uses ordinary file systems rather than
a paging partition. The implementation separates machine-dependent from machine-independent
structure behind a bare validate/invalidate/protect interface --- the independent page size
is a boot-time multiple of the dependent one, 512 bytes on the VAX --- letting the same
machine-independent code sit above VAX page tables and the RT/PC's inverted page table.
Matchmaker, an interface definition language, generates RPC stubs for C, Pascal and
CommonLisp, and user-level network servers extend ports across machines transparently.
As of April 1986 Mach was binary compatible with 4.3BSD and running on assorted VAXen, a
four-processor VAX~11/784 and the IBM RT/PC.

\textbf{Critical judgment.} The gap between the ambition of the design and the weight of the
evidence is very wide, and the paper is candid that ``extensive performance comparisons with
other systems have not yet been done.'' In practice its entire quantitative evaluation is a
single number: on a MicroVAX~II, touching newly allocated memory costs under 0.7\,ms per
1024 bytes against roughly 1.2\,ms for 4.3BSD. That microbenchmark measures zero-fill fault
cost --- exactly what the new VM system was built to improve --- and was chosen by the
implementers; \texttt{fork} is called ``substantially faster'' with no figure attached at all.
Nothing measures IPC latency, although ports and messages are the centrepiece of the design;
nothing measures multiprocessor scaling, although the abstract's opening sentence announces a
``multiprocessor operating system kernel''; threads, the feature the paper leads with, were not yet implemented and were merely
``expected by Summer 1986''; and the caption of the paper's own final figure concedes that
UNIX compatibility still executed in kernel state. The central claim --- that operating
system services can be moved out to user-level servers at acceptable cost --- is therefore
precisely the claim the paper does not test. History then split the verdict. The mechanisms
won comprehensively: separating the address space from the schedulable entity is now
universal, memory-mapped files and external pagers are standard equipment, the
machine-dependent/independent VM split survives as the \texttt{pmap} layer in the BSDs, and
port-like capability handles reappear in Windows object handles and in Fuchsia's channels,
while macOS and iOS run on a Mach derivative to this day. The performance thesis lost: Mach
3.0's user-space UNIX server proved slow, Chen and Bershad later traced the degradation to
memory-system behaviour rather than to message path length, and Liedtke's L4 work argued the
cost was an artefact of implementation rather than of microkernels, cutting IPC by roughly an
order of magnitude. Mach became the cautionary tale that motivated second-generation
microkernels. Two of its assumptions have since expired outright. Network-transparent IPC
treats a local multiprocessor and a LAN as points on one continuum, an equivalence the
distributed systems community abandoned because transparency conceals latency and partial
failure. And copy-on-write as the universal bulk-transfer mechanism assumed page-table
manipulation was cheap relative to copying, which inverted once TLB shootdown across many
cores entered the price. Finally, the extensibility argument is made by \emph{enumerating}
UNIX's accumulated mess rather than by showing that the port model resists the same
accretion --- and it did not: Mach's own VM, with its chains of shadow objects, became
notoriously intricate, and the kernel that carries its name today is large.

\subsection*{Paper 3 --- Baumann et al., ``The Multikernel: A New OS Architecture for Scalable Multicore Systems'' (2009)}
\label{p4:sec:multikernel}

\textbf{Summary.} The Multikernel's premise is that hardware now diversifies faster than
system software can be retuned for it, so that an OS built as a shared-memory kernel with
lock-protected data structures faces an unbounded engineering bill. The authors evidence this
with the ``heroic'' removal of the Windows~7 dispatcher lock, which touched 6000 lines across
58 files, and with the many iterations required by Linux's read-copy-update. They then argue
that the machine has itself become a network --- interconnects route and congest, cores
differ in ISA and performance, and peripherals are already non-coherent --- and propose three
principles: make all inter-core communication explicit, make OS structure hardware-neutral,
and treat state as replicated rather than shared. Their prototype, Barrelfish, factors the
per-core OS instance into a privileged \emph{CPU driver} (event-driven, single-threaded,
non-preemptable, 7135 lines of C for x86-64) and a schedulable user-space \emph{monitor}.
Monitors hold the replicas and run the agreement protocols --- one-phase commit for unmapping
a page, two-phase commit for capability retyping --- over URPC, a user-level message
transport tuned to the cache-coherence protocol, in which the sender writes a cache line
sequentially and the receiver polls its last word. Memory is managed by seL4-style
capabilities with all page-table manipulation performed in user space, and a system knowledge
base holding hardware facts in first-order logic computes topology-aware multicast trees at
run time. The motivating measurement is stark: a single core updates shared state in under 30
cycles, but with 16 cores on the same data each update costs roughly 12{,}000 extra cycles,
and one dedicated RPC server sustains twice the update rate of all 16 shared-memory threads
together. The evaluation then shows a NUMA-aware multicast shootdown protocol holding at
about 3.1k cycles across 32 cores where naive broadcast and unicast climb to roughly 13k;
end-to-end page unmapping at 32 cores costing about 17.5k cycles against roughly 29.5k for
Linux 2.6.26 and 46k for Windows Server 2008~R2; IP loopback at 2154 versus 1823\,Mbit/s with 21
versus 77 data-cache misses per packet; UDP echo over gigabit Ethernet at 951.7 against
951\,Mbit/s; and a static web server at 18{,}697 requests per second against 8924 for
lighttpd on Linux.

\textbf{Critical judgment.} This is much the strongest evaluation of the three and an
unusually honest one --- it states outright that ``it would be wrong to draw any quantitative
conclusions from our large-scale benchmarks,'' concedes that the capability system was ``a
mistake\ldots unnecessarily complex, and no more efficient'' than the per-processor allocators
already in Windows and Linux, calls its network stack ``very much a placeholder,'' and admits
that every test platform was homogeneous x86-64, so the heterogeneity that motivates the
whole argument was never actually exercised. That honesty makes the weak points easier to
locate. The headline numbers are the least defensible: the 2.1$\times$ web-server win pits a
purpose-built user-space lwIP stack against general-purpose lighttpd on Linux with
segmentation and checksum offload enabled, while Barrelfish's own driver supported none of
it, so what is measured is specialisation rather than architecture. The scalability case
rests on global-coordination microbenchmarks, and there the gain comes largely from replacing
serial IPIs with a topology-aware multicast \emph{tree} --- an algorithmic change a
shared-memory kernel could equally adopt --- so the experiment does not isolate shared-nothing
structure as the cause of the improvement. Reading the shootdown and unmap figures against
each other is more revealing still: the raw messaging baseline costs about 3.1k cycles at 32
cores while the complete unmap costs about 17.5k, so more than four-fifths of the operation is
Barrelfish's own monitor LRPC, marshaling and unoptimised thread dispatch. The shape of the
unmap curve makes the same point at the other end. Barrelfish does not start cheaper than the
systems it is compared against --- at two cores it costs about 10.3k cycles against roughly
5.2k for Linux and 5.8k for Windows --- and it stays almost perfectly flat out to five cores
before rising gently. What it buys is slope: from two to thirty-two cores Barrelfish grows
about 1.7$\times$, reaching some 17.5k cycles, where Linux grows 5.7$\times$ to about 29.5k
and Windows 7.9$\times$ to about 46k. So the multikernel did not reduce the constant cost of
the operation; it roughly \emph{doubled} it and bought a much flatter curve in exchange. That
is exactly the trade the authors themselves concede when they note that routing OS invocations
through monitor RPC rather than system calls adds ``a constant overhead on current hardware of
several thousand cycles,'' which is ``constant as the number of cores increases'' --- a
bargain that pays only once the machine is wide enough, and whose break-even against Linux
does not arrive until around twelve or thirteen cores. The messaging table quietly reinforces
the point, showing URPC at 450 cycles of latency against 424 for L4's decade-older IPC, the
real gains again being throughput and a much smaller cache footprint rather than latency. The hardware bet was half right. Cache
coherence did not collapse --- commodity servers remained coherent far past 2009 --- and
mainstream kernels absorbed NUMA and core asymmetry using per-CPU data, RCU and sharded
locks, settling in the middle of the paper's own spectrum of sharing disciplines rather than
at the shared-nothing extreme it advocates. Barrelfish remained a research vehicle. What
survived is the diagnosis and the method: treating the machine as a network, which resurfaced
in kernel-bypass and dataplane systems such as Arrakis and IX, and in the now-routine practice
of designing OS state per-core first. One irony deserves noting --- having adopted the
distributed systems analogy, the paper declines its principal prize, stating that ``the
complete system is a failure unit,'' so it pays the full price of messaging and replication
without buying any tolerance of partial failure.

\subsection*{The Three Together}
\label{p4:sec:together}

Placed side by side, the three papers are three attempts to answer the same question --- what
may the operating system take for granted about the machine? --- and each is undone at a
different layer. Multics assumed that generality could be centralised in one comprehensive
system, and was defeated by the complexity budget it never wrote down. Mach assumed the
boundary could be pushed outward so that services became user-level servers, and was defeated
by the cost of crossing that boundary. The Multikernel assumed shared memory was itself the
wrong primitive, and was overtaken by shared-memory hardware that went on scaling for another
decade. The pattern in their evidence is just as instructive: 1965 offers a design prospectus
with no measurements, 1986 a status report resting on one microbenchmark of the subsystem its
authors had just optimised, and 2009 a broad evaluation that explicitly marks its own limits
--- yet the boldness of the architectural claim does not diminish as the evidence improves,
which is a useful caution when reading any systems paper. Most striking is that all three
have been vindicated in mechanism while being refuted in scope. Multics' paging, protection
rings and device independence are everywhere, but its segmented address space is gone. Mach's
tasks, threads, ports and memory objects are everywhere, but its argument that services could
be evicted from the kernel for free is not. The multikernel's per-core replication now sits
inside Linux, but as an optimisation within a shared-memory kernel rather than as a
replacement for one. The durable contribution of an architecture paper, on this evidence, is
rarely the architecture.

\subsection*{References}
\label{p4:sec:refs}

\begin{enumerate}
  \item[{[1]}] F.~J. Corbat\'o and V.~A. Vyssotsky. ``Introduction and Overview of the
    Multics System.'' In \emph{Proceedings of the AFIPS Fall Joint Computer Conference
    (FJCC)}, vol.~27, part~1, pp.~185--196, 1965.
  \item[{[2]}] M. Accetta, R. Baron, W. Bolosky, D. Golub, R. Rashid, A. Tevanian and
    M. Young. ``Mach: A New Kernel Foundation for UNIX Development.'' In \emph{Proceedings
    of the USENIX Summer Conference}, pp.~93--112, 1986.
  \item[{[3]}] A. Baumann, P. Barham, P.-E. Dagand, T. Harris, R. Isaacs, S. Peter,
    T. Roscoe, A. Sch\"upbach and A. Singhania. ``The Multikernel: A New OS Architecture for
    Scalable Multicore Systems.'' In \emph{Proceedings of the 22nd ACM Symposium on Operating
    Systems Principles (SOSP)}, pp.~29--44, 2009.
\end{enumerate}
